% GENERATED FILE. Edit canonical lessons, then run npm run generate:books.
% canonical-source-hash: fnv1a64:7632f65322254b8d
% canonical-lessons: RU-C468-zhara, RU-C468-dushnyy, RU-C468-dukhota, RU-C468-potet, RU-C468-tayat, RU-R468-review-467-468, RU-R468-review-final-468

\chapter{Heat}
\label{ch:ru-heat}
\hlchaptermodality{468}

\begin{chapteropening}
\textbf{By the end of this chapter:} \emph{I can talk about the heat: hot weather, stuffy air, a close room, sweating and melting.}

Complete the last lesson of chapter 468: I can talk about the heat: hot weather, stuffy air, a close room, sweating and melting.
\end{chapteropening}

\section[\texorpdfstring{\ru{жара} (zhará)}{жара (zhará)}]{\ru{жара} (zhará) — hot weather, heat}

\label{lesson:RU-C468-zhara}



\begin{quote}
Before the new one: say the Russian for to feel the cold, then the Russian for to slip.
\end{quote}

\subsection*{You'll want to know: \ru{жара}}
\textbf{\ru{жара}} (\emph{zhará}) — \textquotedblleft{}hot weather, heat\textquotedblright{}.

\begin{quote}
\ru{В} \ru{июле} \ru{в} \ru{городе} \ru{стояла} \ru{жара}.
\end{quote}

\begin{quote}
\emph{V iyúle v górode stoyála zhará.}
\end{quote}

\begin{quote}
In July the town had a spell of hot weather.
\end{quote}

Feminine. \emph{\ru{В} \ru{жару}} — in hot weather.

\begin{cousinweb}
\emph{\ru{жара}} goes with \emph{\ru{жар}} 'heat'. EtymDB takes \emph{\ru{жара}} from Proto-Slavic \emph{žārĭ} and the PIE root \emph{g\textsuperscript{w}\textsuperscript{h}er-} 'to heat'; the same root is in \emph{\ru{гореть}} 'to burn'.
\end{cousinweb}

\begin{grammarlens}[title={What you've built}]
One more word for telling, explaining and describing at B1.
\end{grammarlens}

\subsection*{Your turn}
\begin{itemize}
\raggedright
  \item \emph{Say it:} \emph{zhará}
  \item \emph{Say it:} \emph{zhará}, once more
  \item \emph{Say it:} a sentence of your own with \emph{zhará}
  \item \emph{Recall:} say the Russian for to feel the cold, then the Russian for to slip, then say \emph{zhará} again
\end{itemize}

\begin{tcolorbox}[breakable,colback=teal!4,colframe=teal!35!black,title={Before you move on}]
What does \ru{жара} mean? (\textquotedblleft{}Hot weather, heat\textquotedblright{}.) Say it once more.
\end{tcolorbox}
\section[\texorpdfstring{\ru{душный} (dúshnyy)}{душный (dúshnyy)}]{\ru{душный} (dúshnyy) — stuffy, close}

\label{lesson:RU-C468-dushnyy}



\begin{quote}
Before the new one: say the Russian for to slip, then the Russian for hot weather.
\end{quote}

\subsection*{You'll want to know: \ru{душный}}
\textbf{\ru{душный}} (\emph{dúshnyy}) — \textquotedblleft{}stuffy, close\textquotedblright{}.

\begin{quote}
\ru{Был} \ru{душный} \ru{летний} \ru{вечер}.
\end{quote}

\begin{quote}
\emph{Byl dúshnyy létniy vécher.}
\end{quote}

\begin{quote}
It was a stuffy summer evening.
\end{quote}

\emph{\ru{Душно}} — it is stuffy.

\begin{cousinweb}
\emph{\ru{душный}} is built on \emph{\ru{дух}} 'breath, air' with \emph{-\ru{ный}}: air that chokes the breath. EtymDB lists \emph{\ru{душный}} without a deeper source.
\end{cousinweb}

\begin{grammarlens}[title={What you've built}]
One more word for telling, explaining and describing at B1.
\end{grammarlens}

\subsection*{Your turn}
\begin{itemize}
\raggedright
  \item \emph{Say it:} \emph{dúshnyy}
  \item \emph{Say it:} \emph{dúshnyy}, once more
  \item \emph{Say it:} a sentence of your own with \emph{dúshnyy}
  \item \emph{Recall:} say the Russian for to slip, then the Russian for hot weather, then say \emph{dúshnyy} again
\end{itemize}

\begin{tcolorbox}[breakable,colback=teal!4,colframe=teal!35!black,title={Before you move on}]
What does \ru{душный} mean? (\textquotedblleft{}Stuffy, close\textquotedblright{}.) Say it once more.
\end{tcolorbox}
\section[\texorpdfstring{\ru{духота} (dukhotá)}{духота (dukhotá)}]{\ru{духота} (dukhotá) — stuffiness, close heat}

\label{lesson:RU-C468-dukhota}



\begin{quote}
Before the new one: say the Russian for hot weather, then the Russian for stuffy.
\end{quote}

\subsection*{You'll want to know: \ru{духота}}
\textbf{\ru{духота}} (\emph{dukhotá}) — \textquotedblleft{}stuffiness, close heat\textquotedblright{}.

\begin{quote}
\ru{В} \ru{вагоне} \ru{была} \ru{ужасная} \ru{духота}.
\end{quote}

\begin{quote}
\emph{V vagóne bylá uzhásnaya dukhotá.}
\end{quote}

\begin{quote}
It was terribly stuffy in the carriage.
\end{quote}

Feminine.

\begin{cousinweb}
\emph{\ru{духота}} is built on \emph{\ru{дух}} 'breath, air' with \emph{-\ru{ота}}: heavy air one can hardly breathe. EtymDB does not list it.
\end{cousinweb}

\begin{grammarlens}[title={What you've built}]
One more word for telling, explaining and describing at B1.
\end{grammarlens}

\subsection*{Your turn}
\begin{itemize}
\raggedright
  \item \emph{Say it:} \emph{dukhotá}
  \item \emph{Say it:} \emph{dukhotá}, once more
  \item \emph{Say it:} a sentence of your own with \emph{dukhotá}
  \item \emph{Recall:} say the Russian for hot weather, then the Russian for stuffy, then say \emph{dukhotá} again
\end{itemize}

\begin{tcolorbox}[breakable,colback=teal!4,colframe=teal!35!black,title={Before you move on}]
What does \ru{духота} mean? (\textquotedblleft{}Stuffiness, close heat\textquotedblright{}.) Say it once more.
\end{tcolorbox}
\section[\texorpdfstring{\ru{потеть} (potét)}{потеть (potét)}]{\ru{потеть} (potét) — to sweat}

\label{lesson:RU-C468-potet}



\begin{quote}
Before the new one: say the Russian for stuffy, then the Russian for stuffiness.
\end{quote}

\subsection*{You'll want to know: \ru{потеть}}
\textbf{\ru{потеть}} (\emph{potét}) — \textquotedblleft{}to sweat\textquotedblright{}.

\begin{quote}
\ru{В} \ru{жару} \ru{я} \ru{сильно} \ru{потею}.
\end{quote}

\begin{quote}
\emph{V zharú ya sílno potéyu.}
\end{quote}

\begin{quote}
In hot weather I sweat a lot.
\end{quote}

Imperfective; perfective \emph{\ru{вспотеть}}. Present \emph{\ru{потею}, \ru{потеешь}}. \emph{\ru{Окна} \ru{потеют}} — the windows mist up.

\begin{cousinweb}
\emph{\ru{потеть}} is built on \emph{\ru{пот}} 'sweat', which is inherited from Proto-Slavic \emph{potŭ} 'sweat' (EtymDB); EtymDB lists \emph{\ru{потеть}} itself without a source.
\end{cousinweb}

\begin{grammarlens}[title={What you've built}]
One more word for telling, explaining and describing at B1.
\end{grammarlens}

\subsection*{Your turn}
\begin{itemize}
\raggedright
  \item \emph{Say it:} \emph{potét}
  \item \emph{Say it:} \emph{potét}, once more
  \item \emph{Say it:} a sentence of your own with \emph{potét}
  \item \emph{Recall:} say the Russian for stuffy, then the Russian for stuffiness, then say \emph{potét} again
\end{itemize}

\begin{tcolorbox}[breakable,colback=teal!4,colframe=teal!35!black,title={Before you move on}]
What does \ru{потеть} mean? (\textquotedblleft{}To sweat\textquotedblright{}.) Say it once more.
\end{tcolorbox}
\section[\texorpdfstring{\ru{таять} (táyat)}{таять (táyat)}]{\ru{таять} (táyat) — to melt, to thaw}

\label{lesson:RU-C468-tayat}



\begin{quote}
Before the new one: say the Russian for stuffiness, then the Russian for to sweat.
\end{quote}

\subsection*{You'll want to know: \ru{таять}}
\textbf{\ru{таять}} (\emph{táyat}) — \textquotedblleft{}to melt, to thaw\textquotedblright{}.

\begin{quote}
\ru{Мороженое} \ru{быстро} \ru{тает} \ru{на} \ru{солнце}.
\end{quote}

\begin{quote}
\emph{Morózhenoye býstro táyet na sólntse.}
\end{quote}

\begin{quote}
The ice cream melts quickly in the sun.
\end{quote}

Imperfective; perfective \emph{\ru{растаять}}. Present \emph{\ru{таю}, \ru{таешь}}.

\begin{cousinweb}
\emph{\ru{таять}} is inherited from Proto-Slavic \emph{tajati} (EtymDB), with the sense 'to melt, to thaw'.
\end{cousinweb}

\begin{grammarlens}[title={What you've built}]
One more word for telling, explaining and describing at B1.
\end{grammarlens}

\subsection*{Your turn}
\begin{itemize}
\raggedright
  \item \emph{Say it:} \emph{táyat}
  \item \emph{Say it:} \emph{táyat}, once more
  \item \emph{Say it:} a sentence of your own with \emph{táyat}
  \item \emph{Recall:} say the Russian for stuffiness, then the Russian for to sweat, then say \emph{táyat} again
\end{itemize}

\begin{tcolorbox}[breakable,colback=teal!4,colframe=teal!35!black,title={Before you move on}]
What does \ru{таять} mean? (\textquotedblleft{}To melt, to thaw\textquotedblright{}.) Say it once more.
\end{tcolorbox}
\section[(dialogue)]{Town, shops, nature and weather — a review}

\label{lesson:RU-R468-review-467-468}



\begin{quote}
Say the Russian for frost and black ice.
\end{quote}

\subsection*{Your turn}
Say these aloud:

\begin{itemize}
\raggedright
  \item frost, black ice, a blizzard, to feel the cold, to slip
  \item hot weather, stuffy, stuffiness, to sweat, to melt
\end{itemize}

\begin{tcolorbox}[breakable,colback=teal!4,colframe=teal!35!black,title={Before you move on}]
Frost? (\textbf{\ru{мороз}}.) Black ice? (\textbf{\ru{гололёд}}.) A blizzard? (\textbf{\ru{метель}}.) To feel the cold? (\textbf{\ru{мёрзнуть}}.) To slip? (\textbf{\ru{поскользнуться}}.) Hot weather? (\textbf{\ru{жара}}.) Stuffy? (\textbf{\ru{душный}}.) Stuffiness? (\textbf{\ru{духота}}.) To sweat? (\textbf{\ru{потеть}}.) To melt? (\textbf{\ru{таять}}.)
\end{tcolorbox}
\section[(dialogue)]{Heat — once more}

\label{lesson:RU-R468-review-final-468}



\begin{quote}
Say the Russian for hot weather and stuffy.
\end{quote}

\subsection*{Your turn}
\begin{itemize}
\raggedright
  \item \emph{Say it:} hot weather, stuffy, stuffiness, to sweat, to melt
\end{itemize}

\begin{tcolorbox}[breakable,colback=teal!4,colframe=teal!35!black,title={Before you move on}]
Hot weather? (\textbf{\ru{жара}}.) Stuffy? (\textbf{\ru{душный}}.) Stuffiness? (\textbf{\ru{духота}}.) To sweat? (\textbf{\ru{потеть}}.) To melt? (\textbf{\ru{таять}}.)
\end{tcolorbox}
